% Fuentes: CLAUDE.md (contexto), EXP-001, DEC-006; xview_ann_train.json.
\section{Objetivo y datos}
La tarea es clasificar recortes de \res{baseline-res}~$\times$~\res{baseline-res} píxeles (RGB) en \res{n-classes}
clases de objetos con una red \emph{feedforward}, eligiendo su arquitectura, el algoritmo de optimización y los
parámetros de entrenamiento, y usando un conjunto de validación para ajustar los hiperparámetros y decidir cuándo
parar. Hay \res{n-images} imágenes etiquetadas, muy desbalanceadas: la clase más frecuente (\res{class-max-name},
\res{class-max-count}) tiene \res{imbalance-ratio} veces más ejemplos que la menos frecuente (\res{class-min-name},
\res{class-min-count}). Se separa un \res{valid-share}\,\% para validación con un \emph{split} estratificado y
semilla fija (\res{n-train} imágenes de entrenamiento y \res{n-valid} de validación), de modo que cada clase conserva
su proporción. El conjunto de test (\res{n-test} imágenes) no tiene etiquetas locales y no se usa para decidir nada.

Por el desbalance, la métrica guía es el \emph{mean recall} (recall medio por clase), cuyo valor por azar es
\res{chance-mean-recall}; predecir siempre la clase mayoritaria da una \emph{accuracy} de \res{majority-acc}. Todas
las cifras de este informe son de validación.
